\label{editorial:prototipo}
\section{Prototipo local para consultar predicciones y explicaciones}
Se implementó un prototipo local de consulta, accesible como archivo autónomo desde un navegador y alimentado por los resultados congelados. Permite elegir el corte temporal, consultar tamaños de población, métricas, intervalos de AP y errores, comparar las familias e inspeccionar la calibración con sus denominadores. La familia principal se distingue de las demás sin permitir sustituirla por resultados de prueba.

La consulta de SHAP presenta magnitudes medias y explicaciones individuales de las cien inscripciones seleccionadas por corte. En cada caso se muestran la probabilidad, el valor de referencia, los valores observados de los indicadores, las contribuciones positivas y negativas y su suma. Las barras parten de cero: el azul representa aumentos de la probabilidad estimada y el verde salvia, reducciones. El signo no expresa una valoración del estudiante ni un efecto causal. Los casos usan una numeración local sin exponer identificadores de estudiantes. La interfaz conserva la regla de alerta congelada; no permite alterar umbrales, reajustar modelos ni producir predicciones nuevas.

Se incorporó también una consulta descriptiva por edad, discapacidad registrada, privación del área de residencia y región, con conteos y porcentajes en los cinco cortes. Esta vista representa la población elegible completa y se distingue explícitamente del conjunto de prueba. Se cotejaron 145 resultados descriptivos con el anexo; esta incorporación no comprende todavía toda la exploración académica, conductual y temporal del proyecto.

La lectura comienza con los retiros detectados y las falsas alertas, continúa con la comparación entre modelos y la calibración, y termina con las explicaciones y el contexto. Se utilizan títulos explicativos, color selectivo, etiquetas directas y tablas desplegables para conservar el detalle. Las dos barras iniciales tienen denominadores distintos: una describe todas las alertas y otra todos los retiros. Los intervalos de incertidumbre permanecen visibles en la comparación de AP.

Las comprobaciones funcionales cubrieron cinco cortes, 500 casos individuales y 20 consultas de contexto mediante una representación mínima de los elementos de la interfaz. No sustituyen una prueba completa de navegador. Los autores confirmaron la apertura de una versión anterior y aportaron capturas; la revisión visual del rediseño final y la evaluación de comprensión con los usuarios previstos continúan pendientes. Se documentó una matriz de cobertura de requisitos y sus limitaciones. El prototipo es académico: no se conecta a un LMS ni ejecuta intervenciones; no demuestra utilidad institucional ni reducción del abandono.
